\chapter{BACKGROUND}

% ================================================
% 1.1 Overview
% ================================================
\section{Overview}

Data quality (DQ) is a critical concern in modern data-intensive systems. As organizations increasingly rely on real-time data streams for decision-making, ensuring the quality of streaming data has become a fundamental challenge. This chapter reviews the foundational concepts, existing frameworks, and relevant datasets that form the basis of this research.

% ================================================
% 1.2 Streaming Data Processing
% ================================================
\section{Streaming Data Processing}

Streaming data processing refers to the continuous ingestion, processing, and analysis of data as it arrives, rather than in predefined batches. Key technologies in this space include Apache Kafka \citep{kafkaintro2024} for distributed event streaming, and Apache Flink \citep{flinkdocs2024} for stateful stream processing with exactly-once semantics.

\subsection{Apache Kafka: Event Streaming Platform}

According to the official Apache Kafka documentation \citep{kafkaintro2024}, Kafka is a distributed event streaming platform that combines three key capabilities: (1) publishing (write) and subscribing to (read) streams of events; (2) storing streams of events durably and reliably; and (3) processing streams of events as they occur or retrospectively. Kafka operates as a distributed system consisting of servers (brokers) and clients communicating via a high-performance TCP network protocol. Kafka is deployed on bare-metal hardware, virtual machines, containers, and in cloud environments.

The Kafka architecture centers on several core concepts \citep{kafkaintro2024}:

\begin{itemize}
    \item \textbf{Events (records/messages):} An event records the fact that ``something happened.'' Conceptually, an event has a key, value, timestamp, and optional metadata headers.
    \item \textbf{Topics:} Events are organized and durably stored in topics---similar to a folder in a filesystem, with events as files. Topics are \textit{multi-producer} and \textit{multi-subscriber}: a topic can have zero, one, or many producers writing events, and zero, one, or many consumers subscribing to events.
    \item \textbf{Partitions:} Topics are partitioned, meaning a topic is spread over a number of ``buckets'' located on different Kafka brokers. This enables both write and read parallelism across multiple brokers simultaneously. Events with the same event key are written to the same partition, guaranteeing that any consumer of a given topic-partition will always read that partition's events in exactly the same order as they were written.
    \item \textbf{Brokers:} Kafka runs as a cluster of servers forming the storage layer. Each broker holds a subset of partitions and serves as the leader for specific partitions. If any server fails, other servers take over their work to ensure continuous operations without data loss.
    \item \textbf{Producers:} Client applications that publish (write) events to Kafka, directly addressing messages to particular partitions.
    \item \textbf{Consumers:} Client applications that subscribe to (read) and process events. Kafka consumers track the maximum offset they have consumed in each partition and can commit offsets to resume from those positions after restart. Consumer groups enable parallel consumption of partitions.
    \item \textbf{Replication:} Every topic can be replicated across brokers for fault tolerance and high availability. A common production setting is a replication factor of 3.
\end{itemize}

Kafka's retention mechanism stores events durably for as long as needed (configurable per topic, e.g., 7 days). Unlike traditional messaging systems, events are not deleted after consumption. This enables replay of historical events, which is essential for backfilling baselines and reprocessing after failures. Kafka's performance remains effectively constant with respect to data size.

The Kafka Streams API \citep{kafkastreams2024} enables implementing stream processing applications and microservices with higher-level functions including transformations, stateful operations (aggregations, joins), windowing, and event-time processing. Kafka Streams partitions data for processing using the same partitioning as Kafka's storage layer, enabling data locality, elasticity, scalability, high performance, and fault tolerance.

\subsection{Apache Flink: Stateful Stream Processing}

Apache Flink \citep{flinkdocs2024} is described as ``a framework and distributed processing engine for stateful computations over unbounded and bounded data streams.'' Flink has been designed to run in all common cluster environments, perform computations at in-memory speed, and scale to any size. Flink integrates with Hadoop YARN, Kubernetes, and can run as a stand-alone cluster.

The Flink architecture is built around four critical concepts \citep{flinkdocs2024}:

\begin{itemize}
    \item \textbf{Continuous processing of streaming data:} Flink applications are composed of streaming dataflows that may be transformed by user-defined operators. These dataflows form directed graphs that start with one or more sources and end in one or more sinks. Applications can consume real-time data from streaming sources (e.g., Kafka, Kinesis) and produce results to various sink systems.
    \item \textbf{Parallel dataflows:} Programs in Flink are inherently parallel and distributed. During execution, a stream has one or more stream partitions, and each operator has one or more operator subtasks. Stream transport between operators can be either \textit{one-to-one} (preserving partitioning and ordering) or \textit{redistributing} (changing partitioning via keyBy, broadcast, or rebalance).
    \item \textbf{Event time processing:} Flink supports using event time timestamps recorded in the data stream, rather than processing time (machine clocks). This enables consistent, accurate analytics regardless of when events are processed, and allows reasoning about when a set of events is complete.
    \item \textbf{Stateful stream processing:} Flink operations can be stateful---how one event is handled can depend on the accumulated effect of all events that came before it. A Flink application running in parallel is effectively a sharded key-value store: each parallel instance of a stateful operator is responsible for handling events for a specific group of keys, with state kept locally for high throughput and low latency.
\end{itemize}

Flink provides exactly-once state consistency through a combination of state snapshots (checkpoints) and stream replay. These snapshots capture the entire state of the distributed pipeline---recording offsets into input queues as well as state throughout the job graph. When a failure occurs, sources are rewound, state is restored, and processing resumes. Flink's asynchronous and incremental checkpointing algorithm ensures minimal impact on processing latencies.

Flink supports both \textit{unbounded streams} (continuous, no defined end, must be continuously processed) and \textit{bounded streams} (batch processing with a defined start and end). Bounded streams are internally processed by algorithms and data structures specifically designed for fixed-size datasets, yielding excellent performance. This makes Flink a unified batch-stream processing engine.

Flink's JobManager coordinates the execution of a Flink job, managing task scheduling, checkpoint coordination, and failure recovery. TaskManagers are worker nodes that execute tasks in parallel within processing slots. Each TaskManager provides a fixed number of processing slots (controlled by taskmanager.numberOfTaskSlots). State is maintained in memory or, if state size exceeds available memory, in access-efficient on-disk data structures (RocksDB state backend). Applications can leverage virtually unlimited amounts of CPUs, main memory, disk, and network I/O, with Flink reported to process multiple trillions of events per day and maintain multiple terabytes of state in production deployments.

Flink's windowing mechanism enables aggregation over finite sets of events. Windows can be defined by time (tumbling, sliding, session) or count, and can be assigned based on event time or processing time.

The distinction between batch and streaming DQ is fundamental. In batch processing, all records are available before validation, enabling complete statistical analysis. In streaming, records arrive incrementally, requiring online algorithms that update statistics incrementally and validate records with minimal latency.


% ================================================
% 1.3 MLOps: Machine Learning Operations
% ================================================
\section{MLOps: Machine Learning Operations}

MLOps (Machine Learning Operations) is an ML engineering culture and practice that aims at unifying ML system development (Dev) and ML system operation (Ops) \citep{googlemlops2024}. Practicing MLOps means advocating for automation and monitoring at all steps of ML system construction, including integration, testing, releasing, deployment, and infrastructure management. Only a small fraction of a real-world ML system is composed of the ML code -- the required surrounding elements are vast and complex.

The essential MLOps system components are:

\begin{itemize}
    \item \textbf{Configuration:} Model hyperparameters, feature engineering parameters, pipeline configurations.
    \item \textbf{Automation:} CI/CD pipelines for training and serving.
    \item \textbf{Data Collection:} Continuous ingestion of training and validation data.
    \item \textbf{Data Verification:} Automated validation of data quality and schema.
    \item \textbf{Testing and Debugging:} Unit tests for code, validation tests for data and models.
    \item \textbf{Resource Management:} Compute infrastructure for training and inference.
    \item \textbf{Model Analysis:} Evaluation metrics, fairness metrics, behavioral testing.
    \item \textbf{Process and Metadata Management:} Experiment tracking, model versioning, lineage.
    \item \textbf{Serving Infrastructure:} Model deployment, A/B testing, canary rollouts.
    \item \textbf{Monitoring:} Production model performance, data drift, concept drift detection.
\end{itemize}

MLflow \citep{mlflow2024} is the industry-standard open-source platform for MLOps, providing experiment tracking, model versioning, model registry, and serving capabilities for ML system lifecycle management.

\subsection{MLOps Maturity Levels}

The MLOps maturity model \citep{googlemlops2024} defines three progressive levels:

\begin{itemize}
    \item \textbf{Level 0 (Manual):} Entire ML pipeline is manual. Data scientists train models in notebooks, hand off to engineers for deployment. No CI/CD, no continuous training (CT). Suitable for systems where the model rarely changes.
    \item \textbf{Level 1 (Automated ML Pipeline):} The ML pipeline is automated so that it continuously trains models in production. CT is enabled via automated data and model validation, triggered by scheduled jobs or data/model changes. Still typically manual for deploying model updates.
    \item \textbf{Level 2 (Automated ML with CI/CD):} The ML pipeline is automated with a robust CI/CD system. Data scientists can rapidly explore new features, algorithms, and model architectures. The system automatically builds, tests, and deploys new pipeline components. This level enables rapid iteration and production-grade reliability.
\end{itemize}

CA-DQStream targets MLOps Level 2, as detailed in Chapter 4.

% ================================================
% 1.4 DAMA-DMBOK Data Quality Dimensions
% ================================================
\section{DAMA-DMBOK Data Quality Dimensions}

The Data Management Association (DAMA) and the Data Management Body of Knowledge (DMBOK) define six core data quality dimensions that guide DQ assessment \citep{dama2017dmbok}:

\begin{enumerate}
    \item \textbf{Completeness:} The degree to which data is not missing. Measured by the proportion of NULL or blank values in required fields.
    \item \textbf{Validity:} Data conforms to defined syntax and semantic rules. For example, fare amounts must be non-negative.
    \item \textbf{Accuracy:} The degree to which data reflects the real-world phenomenon it describes. Detected through outlier analysis and cross-validation.
    \item \textbf{Consistency:} Freedom from contradiction within or across datasets. Cross-feature consistency rules detect values that contradict each other (e.g., cash payments should have zero tip).
    \item \textbf{Timeliness:} Data is current enough for its intended use. Measured by event-time vs. processing-time lag.
    \item \textbf{Uniqueness:} No duplicate records exist. Detected through key-based deduplication.
\end{enumerate}

CA-DQStream addresses all six dimensions through its layered architecture, as detailed in Chapter 4.

% ================================================
\section{Chapter Summary}
% ================================================

This chapter reviewed the foundational concepts for this thesis. Streaming data processing with Apache Kafka and Apache Flink provides the distributed infrastructure for real-time DQ monitoring. The MLOps maturity model establishes the operational context in which CA-DQStream operates (Level 2 automation with drift-triggered retraining). The DAMA-DMBOK framework defines six core DQ dimensions, all of which CA-DQStream addresses through its layered architecture. The next chapter surveys the existing literature to identify specific gaps that this thesis addresses.
